% GENERATED FILE. Edit canonical lessons, then run npm run generate:books.
% canonical-source-hash: fnv1a64:b5c5175d669163c7
% canonical-lessons: KA-C522-kammara, KA-C522-akkasaliga, KA-C522-nekara, KA-C522-silpi, KA-C522-kettu

\chapter{Craftspeople}
\label{ch:ka-craftspeople}
\hlchaptermodality{522}

\begin{chapteropening}
\textbf{By the end of this chapter:} \emph{I can say a blacksmith, a goldsmith, a weaver, a sculptor and to carve.}

Complete the last lesson of chapter 522: I can say a blacksmith, a goldsmith, a weaver, a sculptor and to carve.
\end{chapteropening}

\section[\texorpdfstring{\kn{ಕಮ್ಮಾರ} (kammāra)}{ಕಮ್ಮಾರ (kammāra)}]{\kn{ಕಮ್ಮಾರ} (kammāra) — a blacksmith}

\label{lesson:KA-C522-kammara}



\begin{quote}
Before the new one: say the Kannada for hard work, then the Kannada for rivalry.
\end{quote}

\subsection*{You'll want to know: \kn{ಕಮ್ಮಾರ}}
\textbf{\kn{ಕಮ್ಮಾರ}} — \emph{kammāra} — \textquotedblleft{}a blacksmith\textquotedblright{}.

\begin{quote}
\kn{ಕಮ್ಮಾರ} \kn{ಬೆಂಕಿಯಲ್ಲಿ} \kn{ಕಬ್ಬಿಣವನ್ನು} \kn{ಕಾಯಿಸಿದನು}.
\end{quote}

\begin{quote}
\emph{kammāra beṅkiyalli kabbiṇavannu kāyisidanu.}
\end{quote}

\begin{quote}
The blacksmith heated the iron in the fire.
\end{quote}

Plural: \kn{ಕಮ್ಮಾರರು}, \emph{kammāraru}.

His forge is the \kn{ಕುಲುಮೆ}.

\begin{cousinweb}
\textbf{\kn{ಕಮ್ಮಾರ}} (\emph{kammāra}) goes back to Sanskrit \emph{karmāra}, \textquotedblleft{}blacksmith\textquotedblright{} (CDIAL 2898), through a Prakrit form in which \emph{rm} became \emph{mm}.
\end{cousinweb}

\begin{grammarlens}[title={What you've built}]
One more word for describing what you lived and felt.
\end{grammarlens}

\subsection*{Your turn}
\begin{itemize}
\raggedright
  \item \emph{Say it:} \emph{kammāra}
  \item \emph{Say it:} \emph{kammāra}, once more
  \item \emph{Say it:} \emph{kammāra beṅkiyalli kabbiṇavannu kāyisidanu.} — The blacksmith heated the iron in the fire.
  \item \emph{Recall:} say the Kannada for hard work, then the Kannada for rivalry, then say \emph{kammāra} again
\end{itemize}

\begin{tcolorbox}[breakable,colback=teal!4,colframe=teal!35!black,title={Before you move on}]
What does \kn{ಕಮ್ಮಾರ} mean? (\textquotedblleft{}A blacksmith\textquotedblright{}.) Say it once more.
\end{tcolorbox}
\section[\texorpdfstring{\kn{ಅಕ್ಕಸಾಲಿಗ} (akkasāliga)}{ಅಕ್ಕಸಾಲಿಗ (akkasāliga)}]{\kn{ಅಕ್ಕಸಾಲಿಗ} (akkasāliga) — a goldsmith}

\label{lesson:KA-C522-akkasaliga}



\begin{quote}
Before the new one: say the Kannada for rivalry, then the Kannada for a blacksmith.
\end{quote}

\subsection*{You'll want to know: \kn{ಅಕ್ಕಸಾಲಿಗ}}
\textbf{\kn{ಅಕ್ಕಸಾಲಿಗ}} — \emph{akkasāliga} — \textquotedblleft{}a goldsmith\textquotedblright{}.

\begin{quote}
\kn{ಅಕ್ಕಸಾಲಿಗ} \kn{ಮದುವೆಗೆ} \kn{ಚಿನ್ನದ} \kn{ಸರ} \kn{ಮಾಡಿದನು}.
\end{quote}

\begin{quote}
\emph{akkasāliga maduvege cinnada sara māḍidanu.}
\end{quote}

\begin{quote}
The goldsmith made a gold necklace for the wedding.
\end{quote}

Plural: \kn{ಅಕ್ಕಸಾಲಿಗರು}, \emph{akkasāligaru}.

The ending -\kn{ಇಗ} marks the worker, as in many trade names.

\begin{cousinweb}
\textbf{\kn{ಅಕ್ಕಸಾಲಿಗ}} (\emph{akkasāliga}) is the Kannada name of the goldsmith; its origin is not settled.
\end{cousinweb}

\begin{grammarlens}[title={What you've built}]
One more word for describing what you lived and felt.
\end{grammarlens}

\subsection*{Your turn}
\begin{itemize}
\raggedright
  \item \emph{Say it:} \emph{akkasāliga}
  \item \emph{Say it:} \emph{akkasāliga}, once more
  \item \emph{Say it:} \emph{akkasāliga maduvege cinnada sara māḍidanu.} — The goldsmith made a gold necklace for the wedding.
  \item \emph{Recall:} say the Kannada for rivalry, then the Kannada for a blacksmith, then say \emph{akkasāliga} again
\end{itemize}

\begin{tcolorbox}[breakable,colback=teal!4,colframe=teal!35!black,title={Before you move on}]
What does \kn{ಅಕ್ಕಸಾಲಿಗ} mean? (\textquotedblleft{}A goldsmith\textquotedblright{}.) Say it once more.
\end{tcolorbox}
\section[\texorpdfstring{\kn{ನೇಕಾರ} (nēkāra)}{ನೇಕಾರ (nēkāra)}]{\kn{ನೇಕಾರ} (nēkāra) — a weaver}

\label{lesson:KA-C522-nekara}



\begin{quote}
Before the new one: say the Kannada for a blacksmith, then the Kannada for a goldsmith.
\end{quote}

\subsection*{You'll want to know: \kn{ನೇಕಾರ}}
\textbf{\kn{ನೇಕಾರ}} — \emph{nēkāra} — \textquotedblleft{}a weaver\textquotedblright{}.

\begin{quote}
\kn{ನೇಕಾರರು} \kn{ಕೈಮಗ್ಗದಲ್ಲಿ} \kn{ರೇಷ್ಮೆ} \kn{ಸೀರೆ} \kn{ನೇಯುತ್ತಾರೆ}.
\end{quote}

\begin{quote}
\emph{nēkāraru kaimaggadalli rēṣme sīre nēyuttāre.}
\end{quote}

\begin{quote}
Weavers weave silk saris on hand looms.
\end{quote}

Plural: \kn{ನೇಕಾರರು}, \emph{nēkāraru}.

The verb is \kn{ನೇಯು}, to weave.

\begin{cousinweb}
\textbf{\kn{ನೇಕಾರ}} (\emph{nēkāra}) is native Dravidian: DEDR 3745 lists Kannada \emph{nēy}, \emph{nēyu}, \textquotedblleft{}to weave\textquotedblright{}, and agent nouns such as \emph{nēyikāṟa}, \textquotedblleft{}a weaver\textquotedblright{}.
\end{cousinweb}

\begin{grammarlens}[title={What you've built}]
One more word for describing what you lived and felt.
\end{grammarlens}

\subsection*{Your turn}
\begin{itemize}
\raggedright
  \item \emph{Say it:} \emph{nēkāra}
  \item \emph{Say it:} \emph{nēkāra}, once more
  \item \emph{Say it:} \emph{nēkāraru kaimaggadalli rēṣme sīre nēyuttāre.} — Weavers weave silk saris on hand looms.
  \item \emph{Recall:} say the Kannada for a blacksmith, then the Kannada for a goldsmith, then say \emph{nēkāra} again
\end{itemize}

\begin{tcolorbox}[breakable,colback=teal!4,colframe=teal!35!black,title={Before you move on}]
What does \kn{ನೇಕಾರ} mean? (\textquotedblleft{}A weaver\textquotedblright{}.) Say it once more.
\end{tcolorbox}
\section[\texorpdfstring{\kn{ಶಿಲ್ಪಿ} (śilpi)}{ಶಿಲ್ಪಿ (śilpi)}]{\kn{ಶಿಲ್ಪಿ} (śilpi) — a sculptor}

\label{lesson:KA-C522-silpi}



\begin{quote}
Before the new one: say the Kannada for a goldsmith, then the Kannada for a weaver.
\end{quote}

\subsection*{You'll want to know: \kn{ಶಿಲ್ಪಿ}}
\textbf{\kn{ಶಿಲ್ಪಿ}} — \emph{śilpi} — \textquotedblleft{}a sculptor\textquotedblright{}.

\begin{quote}
\kn{ಶಿಲ್ಪಿ} \kn{ಕಲ್ಲಿನಲ್ಲಿ} \kn{ಸುಂದರ} \kn{ಆನೆಯನ್ನು} \kn{ಮಾಡಿದನು}.
\end{quote}

\begin{quote}
\emph{śilpi kallinalli sundara āneyannu māḍidanu.}
\end{quote}

\begin{quote}
The sculptor made a beautiful elephant in stone.
\end{quote}

Plural: \kn{ಶಿಲ್ಪಿಗಳು}, \emph{śilpigaḷu}.

\kn{ಶಿಲ್ಪ} is a sculpture.

\begin{cousinweb}
\textbf{\kn{ಶಿಲ್ಪಿ}} (\emph{śilpi}) is borrowed from Sanskrit \emph{śilpin}, \textquotedblleft{}skilled in art; an artificer\textquotedblright{} (CDIAL 12471), built on \emph{śilpa}, \textquotedblleft{}artistic work\textquotedblright{}.
\end{cousinweb}

\begin{grammarlens}[title={What you've built}]
One more word for describing what you lived and felt.
\end{grammarlens}

\subsection*{Your turn}
\begin{itemize}
\raggedright
  \item \emph{Say it:} \emph{śilpi}
  \item \emph{Say it:} \emph{śilpi}, once more
  \item \emph{Say it:} \emph{śilpi kallinalli sundara āneyannu māḍidanu.} — The sculptor made a beautiful elephant in stone.
  \item \emph{Recall:} say the Kannada for a goldsmith, then the Kannada for a weaver, then say \emph{śilpi} again
\end{itemize}

\begin{tcolorbox}[breakable,colback=teal!4,colframe=teal!35!black,title={Before you move on}]
What does \kn{ಶಿಲ್ಪಿ} mean? (\textquotedblleft{}A sculptor\textquotedblright{}.) Say it once more.
\end{tcolorbox}
\section[\texorpdfstring{\kn{ಕೆತ್ತು} (kettu)}{ಕೆತ್ತು (kettu)}]{\kn{ಕೆತ್ತು} (kettu) — to carve; to engrave}

\label{lesson:KA-C522-kettu}



\begin{quote}
Before the new one: say the Kannada for a weaver, then the Kannada for a sculptor.
\end{quote}

\subsection*{You'll want to know: \kn{ಕೆತ್ತು}}
\textbf{\kn{ಕೆತ್ತು}} — \emph{kettu} — \textquotedblleft{}to carve; to engrave\textquotedblright{}.

\begin{quote}
\kn{ಬಡಗಿ} \kn{ಮರದ} \kn{ಬಾಗಿಲಿನ} \kn{ಮೇಲೆ} \kn{ಹೂವುಗಳನ್ನು} \kn{ಕೆತ್ತಿದನು}.
\end{quote}

\begin{quote}
\emph{baḍagi marada bāgilina mēle hūvugaḷannu kettidanu.}
\end{quote}

\begin{quote}
The carpenter carved flowers on the wooden door.
\end{quote}

Past (I …): \kn{ಕೆತ್ತಿದೆ}, \emph{kettide}.

\kn{ಕೆತ್ತನೆ}, a carving.

\begin{cousinweb}
\textbf{\kn{ಕೆತ್ತು}} (\emph{kettu}) is native Dravidian: DEDR 1953 lists Kannada \emph{kettu}, \textquotedblleft{}to pare, chip, engrave, carve\textquotedblright{}.
\end{cousinweb}

\begin{grammarlens}[title={What you've built}]
One more word for describing what you lived and felt.
\end{grammarlens}

\subsection*{Your turn}
\begin{itemize}
\raggedright
  \item \emph{Say it:} \emph{kettu}
  \item \emph{Say it:} \emph{kettu}, once more
  \item \emph{Say it:} \emph{baḍagi marada bāgilina mēle hūvugaḷannu kettidanu.} — The carpenter carved flowers on the wooden door.
  \item \emph{Recall:} say the Kannada for a weaver, then the Kannada for a sculptor, then say \emph{kettu} again
\end{itemize}

\begin{tcolorbox}[breakable,colback=teal!4,colframe=teal!35!black,title={Before you move on}]
What does \kn{ಕೆತ್ತು} mean? (\textquotedblleft{}To carve; to engrave\textquotedblright{}.) Say it once more.
\end{tcolorbox}
